%! Linear Regression — Unit 1, Lesson 1.5 — Student Packet Cover Sheet
\documentclass[10pt]{article}
\usepackage{algebra2-article}
\usepackage{algebra2-boxes}

\begin{document}

% Full-bleed forest banner
\begin{tikzpicture}[remember picture, overlay]
  \fill[forest] (current page.north west)
    rectangle ([yshift=-0.9in]current page.north east);
\end{tikzpicture}
\vspace{-0.8in}

\begin{center}
  {\color{white}\textbf{\LARGE Algebra 2}}\\[4pt]
  {\color{forestmid}\large\textbf{Unit 1 \quad Linear Functions}}\\[6pt]
  {\color{white}\textbf{\large Lesson 1.5 \quad Linear Regression}}
\end{center}

\vspace{0.22in}
\namedateperiod
\vspace{0.18in}

\begin{learningtargetbox}
\small
\begin{enumerate}[leftmargin=1.5em, itemsep=5pt]
  \item \ldots{} read a \textbf{scatterplot} of bivariate data and describe its
        \textbf{association} by \emph{direction}, \emph{form}, and \emph{strength}.
  \item \ldots{} interpret the \textbf{correlation coefficient} $r$ --- its \emph{sign} is the
        direction and its \emph{size} is the strength --- and match an $r$-value to a plot.
  \item \ldots{} read a \textbf{line of best fit} $\hat y=ax+b$ from technology, interpret its
        \textbf{slope} and \textbf{$y$-intercept} in context, and \textbf{predict} with it.
  \item \ldots{} tell \textbf{interpolation} from \textbf{extrapolation}, judge whether a
        prediction is reasonable, and explain why \textbf{correlation is not causation}.
\end{enumerate}
\end{learningtargetbox}

\vspace{0.14in}

\begin{tocbox}
\small
\renewcommand{\arraystretch}{1.5}
\begin{tabularx}{\linewidth}{@{} c l X r @{}}
\rowcolor{forestbg}
\textbf{\#} & \textbf{Component} & \textbf{Description} & \textbf{Score} \\
\midrule
1 & Warm-Up        & Spiral review: reading points off a grid, predicting with a rule, and
                     slope as direction                                                        & \blank{1.2cm} \\
2 & Guided Notes   & Scatterplots and association; the correlation coefficient; the line of
                     best fit and what its numbers mean; predicting and its limits --- ending in
                     \emph{Guided Practice} and \emph{Individual Practice}                     & \blank{1.2cm} \\
3 & Homework       & Mixed practice: matching $r$ to plots, a direction-vs-strength pair, a
                     model from a data table, the special values of $r$, a causation critique
                     --- \emph{due next class}                                                 & \blank{1.2cm} \\
\midrule
  & \multicolumn{2}{r}{\textbf{Total}} & \blank{1.2cm} \\
\end{tabularx}
\end{tocbox}

\vspace{0.10in}

\begin{remindbox}
\small Real data rarely lands on one line, so we summarize the \emph{trend}. A
\textbf{scatterplot}'s association has a \emph{direction} (positive or negative), a \emph{form}
(does a line fit?), and a \emph{strength} (how tightly?). The \textbf{correlation coefficient}
$r$ packs two of these into one number between $-1$ and $1$: its \textbf{sign} is the direction
and its \textbf{size} $|r|$ is the strength --- so $r=-0.9$ is a \emph{strong} fit. The
\textbf{line of best fit} $\hat y=ax+b$ comes from technology; its slope is a \emph{per-unit
rate} and its intercept a \emph{baseline}, both with units. Predicting \emph{inside} the data
(\textbf{interpolation}) is safe; \emph{outside} it (\textbf{extrapolation}) the trend may break.
And a strong $r$ never proves that one variable \emph{causes} the other.
\end{remindbox}

\end{document}
